\section{Variational and multisymplectic geometry of the auxiliary field}
\label{sec:variational-multisymplectic}

The linearizing potential does more than simplify the velocity-field
equation. It places the regular constant-$\ellzero$ sheets inside an exact
covariant Hamiltonian system, identifies the lower density root as a
degeneracy of the Legendre map, and turns the Green identities into a
multisymplectic conservation law. This structure is useful only if its
logical level is kept clear. It belongs to the auxiliary meridional equation;
it is not, by itself, the physical phase space or energy of the complete
Einstein--dust system. We use the standard variational-bicomplex,
multisymplectic, and covariant-polymomentum conventions of
Refs.~\cite{Anderson1992VariationalBicomplex,MarsdenPatrickShkoller1998CMP,
Kanatchikov1998RepMathPhys}.

Throughout this section, $D$ is an oriented, connected, bounded domain with piecewise
$C^2$ boundary whose closure is contained in $\{r>0\}$, and
$d^2x:=dr\wedge dz$. Repeated meridional indices $i,j\in\{r,z\}$ are summed
with the Euclidean metric. The outward unit normal and arclength measure on
$\partial D$ are denoted by $n$ and $ds$, respectively.

\subsection{Weighted action and its singular pullback}

Recall the auxiliary operator
\begin{equation}
\mathcal L\mathcal F
:=\mathcal F_{rr}-\frac{1}{r}\mathcal F_r+\mathcal F_{zz}
=r\,\partial_i\!\left(\frac{\partial_i\mathcal F}{r}\right).
\label{eq:ms-weighted-operator}
\end{equation}
Its natural action is
\begin{equation}
\mathcal A_D[\mathcal F]
:=\frac12\int_D\frac{|\nabla\mathcal F|^2}{r}\,dr\,dz.
\label{eq:ms-weighted-action}
\end{equation}

\begin{proposition}[Weighted variational principle]
\label{prop:ms-weighted-variation}
Let $\mathcal F\in H^1(D)$, and assume sufficient regularity whenever a
classical trace or derivative is used. For every smooth variation $f$,
\begin{equation}
\delta\mathcal A_D[\mathcal F](f)
=-\int_D\frac{\mathcal L\mathcal F}{r}f\,dr\,dz
+\int_{\partial D}\frac{\partial_n\mathcal F}{r}f\,ds.
\label{eq:ms-first-variation}
\end{equation}
Consequently, the critical points for fixed Dirichlet data are precisely the
weak solutions of $\mathcal L\mathcal F=0$. The action is strictly convex on
each affine Dirichlet class and is strictly convex modulo constants when no
Dirichlet trace is fixed.
\end{proposition}

\begin{proof}
Equation~\eqref{eq:ms-first-variation} is the weighted divergence theorem.
The second variation is $\int_D|\nabla f|^2/r\,dr\,dz$, which proves the
convexity statements.
\end{proof}

We next pull back the action through the constant-$\ellzero$ linearizing
map. On $\eta>0$, set
\begin{align}
\Phi_{\ellzero}(r,\eta)
&:=2\eta+\ellzero r^2\log\!\left(\frac{\eta}{\eta_*}\right)
-\frac{\ellzero}{2}\eta^2,
\qquad \eta_*>0,
\label{eq:ms-Phi-definition}\\
A(r,\eta)
&:=\partial_\eta\Phi_{\ellzero}
=2+\frac{\ellzero r^2}{\eta}-\ellzero\eta,
\qquad
B(r,\eta)
:=\left.\partial_r\Phi_{\ellzero}\right|_\eta
=2\ellzero r\log\!\left(\frac{\eta}{\eta_*}\right).
\label{eq:ms-A-B-definition}
\end{align}
Thus
\begin{equation}
\mathcal F_r=B+A\eta_r,
\qquad
\mathcal F_z=A\eta_z.
\label{eq:ms-Phi-total-derivatives}
\end{equation}

\begin{theorem}[Variational equivalence and fold degeneracy]
\label{thm:ms-pullback-fold}
The pullback of Eq.~\eqref{eq:ms-weighted-action} is the first-order action
with density
\begin{equation}
L_\eta
=\frac{(B+A\eta_r)^2+A^2\eta_z^2}{2r}.
\label{eq:ms-pullback-density}
\end{equation}
Its Euler--Lagrange derivative and fibre Hessian are
\begin{align}
\frac{\delta L_\eta}{\delta\eta}
&=-\frac{A}{r}\mathcal L\mathcal F,
\label{eq:ms-pullback-EL}\\
\frac{\partial^2L_\eta}{\partial\eta_i\partial\eta_j}
&=\frac{A^2}{r}\delta_{ij}.
\label{eq:ms-pullback-Hessian}
\end{align}
Since $\mathcal L\Phi_{\ellzero}(r,rv)=r\EE[v]$, the pulled-back
Euler--Lagrange equation is exactly $A\EE[v]=0$. It is equivalent to the VFE
on every sheet on which $A\neq0$, but not on the fold $A=0$.

More precisely, for $\eta=rv$ with $0<v<1$,
\begin{equation}
A=0
\quad\Longleftrightarrow\quad
r\ellzero=-\frac{2v}{1-v^2}.
\label{eq:ms-fold-density-root}
\end{equation}
If $\ellzero\neq0$, then at such a point
\begin{equation}
\partial_{\eta\eta}\Phi_{\ellzero}
=-\ellzero\left(1+\frac{1}{v^2}\right)\neq0.
\label{eq:ms-fold-second-derivative}
\end{equation}
Hence the map $(r,\eta)\mapsto(r,\mathcal F)$ has a nondegenerate fold.
The Legendre rank of Eq.~\eqref{eq:ms-pullback-density} is two for $A\neq0$
and zero for $A=0$.
\end{theorem}

\begin{proof}
The first identity is the Euler--Lagrange chain rule for a base-preserving
point transformation: $\delta\mathcal F=A\,\delta\eta$. Direct
differentiation of Eq.~\eqref{eq:ms-pullback-density} gives
Eq.~\eqref{eq:ms-pullback-Hessian}. Equation~\eqref{eq:ms-fold-density-root}
follows by substituting $\eta=rv$ into $A=0$, and
Eq.~\eqref{eq:ms-fold-second-derivative} follows by differentiating $A$ at
fixed $r$.
\end{proof}

\begin{remark}[What the singular action does not prove]
\label{rem:ms-fold-caution}
The factor $A$ in Eq.~\eqref{eq:ms-pullback-EL} creates a spurious
variational stratum: a smooth section whose image lies in $A=0$ is stationary
for the pulled-back action even when $\mathcal L\mathcal F\neq0$. Therefore,
one must impose the inverse-sheet condition or return to the original VFE at
the fold. The rank loss is not a gauge symmetry and does not, by itself,
exclude a specially compatible smooth crossing. Since the rank is not
constant across the fold, a conventional unstratified Dirac--Bergmann
analysis is inapplicable there.
\end{remark}

\subsection{De Donder--Weyl Hamiltonian and Poincar\'e--Cartan form}

Let $Y=D\times\mathbb R_{\mathcal F}\to D$ be the auxiliary configuration
bundle. Its first jet has coordinates
$(r,z,\mathcal F,\mathcal F_r,\mathcal F_z)$, while the restricted
multiphase bundle has coordinates $(r,z,\mathcal F,p^r,p^z)$. With
$d^1x_r=dz$ and $d^1x_z=-dr$, define
\begin{equation}
p^r:=\frac{\mathcal F_r}{r},
\qquad
p^z:=\frac{\mathcal F_z}{r},
\qquad
H_{\mathcal F}:=\frac r2\bigl[(p^r)^2+(p^z)^2\bigr].
\label{eq:ms-F-polymomenta-Hamiltonian}
\end{equation}
Our sign convention for the Hamiltonian Poincar\'e--Cartan and
multisymplectic forms is
\begin{align}
\Theta_{\mathcal F}
&:=p^r\,d\mathcal F\wedge dz
-p^z\,d\mathcal F\wedge dr
-H_{\mathcal F}\,dr\wedge dz,
\label{eq:ms-F-Cartan-form}\\
\Omega_{\mathcal F}&:=-d\Theta_{\mathcal F}.
\label{eq:ms-F-multisymplectic-form}
\end{align}

\begin{theorem}[Exact covariant canonical transformation]
\label{thm:ms-covariant-canonical-map}
The De Donder--Weyl equations determined by
Eqs.~\eqref{eq:ms-F-polymomenta-Hamiltonian}--\eqref{eq:ms-F-multisymplectic-form}
are
\begin{equation}
\mathcal F_r=rp^r,
\qquad
\mathcal F_z=rp^z,
\qquad
\partial_rp^r+\partial_zp^z=0,
\label{eq:ms-F-DDW-equations}
\end{equation}
and are equivalent to $\mathcal L\mathcal F=0$.

On $A\neq0$, the $\eta$ polymomenta and Hamiltonian are
\begin{align}
\pi^r&=\frac{A}{r}(B+A\eta_r)=Ap^r,
&
\pi^z&=\frac{A^2}{r}\eta_z=Ap^z,
\label{eq:ms-eta-polymomenta}\\
H_\eta
&=\frac{r}{2A^2}\bigl[(\pi^r)^2+(\pi^z)^2\bigr]
-\frac BA\pi^r.
\label{eq:ms-eta-Hamiltonian}
\end{align}
The map
\begin{equation}
\Psi:(r,z,\eta,\pi^r,\pi^z)
\longmapsto
\left(r,z,\Phi_{\ellzero}(r,\eta),\frac{\pi^r}{A},
\frac{\pi^z}{A}\right)
\label{eq:ms-canonical-map}
\end{equation}
satisfies
\begin{equation}
\Psi^*\Theta_{\mathcal F}=\Theta_\eta,
\qquad
\Psi^*\Omega_{\mathcal F}=\Omega_\eta,
\label{eq:ms-canonical-pullback}
\end{equation}
where
\begin{equation}
\Theta_\eta
:=\pi^r\,d\eta\wedge dz
-\pi^z\,d\eta\wedge dr
-H_\eta\,dr\wedge dz.
\label{eq:ms-eta-Cartan-form}
\end{equation}
Thus the point linearization is an exact local multisymplectic
transformation on every invertible sheet.
\end{theorem}

\begin{proof}
The first two equations in Eq.~\eqref{eq:ms-F-DDW-equations} are the
polymomentum definitions, and the third is the weighted field equation.
Solving Eq.~\eqref{eq:ms-eta-polymomenta} for the first derivatives gives
\begin{equation}
\eta_r=\frac{r\pi^r}{A^2}-\frac BA,
\qquad
\eta_z=\frac{r\pi^z}{A^2},
\end{equation}
from which Eq.~\eqref{eq:ms-eta-Hamiltonian} follows by Legendre transform.
Finally, $d\mathcal F=A\,d\eta+B\,dr$. Substitution into
Eq.~\eqref{eq:ms-F-Cartan-form} gives Eq.~\eqref{eq:ms-eta-Cartan-form}
coefficient by coefficient. Exterior differentiation proves the second
identity in Eq.~\eqref{eq:ms-canonical-pullback}.
\end{proof}

At $A=0$, the map in Eq.~\eqref{eq:ms-canonical-map} ceases to exist, the
velocity Hessian vanishes, and the fibre directions of the pulled-back
Cartan form acquire a kernel. The appropriate object is therefore a
stratified premultisymplectic system, not a single smooth hyperregular
Hamiltonian theory crossing the density fold.

\subsection{Multisymplectic conservation and Noether currents}

Let $f$ and $g$ be two solutions of the linearized equation, which here is
again $\mathcal Lh=0$. Define the one-form
\begin{equation}
\boldsymbol\omega(f,g)
:=\frac{f g_r-g f_r}{r}\,dz
-\frac{f g_z-g f_z}{r}\,dr.
\label{eq:ms-current-form}
\end{equation}
Direct differentiation gives the off-shell identity
\begin{equation}
d\boldsymbol\omega(f,g)
=\frac{f\mathcal Lg-g\mathcal Lf}{r}\,dr\wedge dz.
\label{eq:ms-form-formula-offshell}
\end{equation}
Consequently,
\begin{equation}
d\boldsymbol\omega(f,g)=0
\label{eq:ms-form-formula}
\end{equation}
for every pair of Jacobi fields. Equation~\eqref{eq:ms-form-formula} is the
multisymplectic form formula for the present theory. It is also the local
Green identity and hence makes precise why the same conserved quantity
controls solution transformations, boundary data, and numerical residuals.

Several exact Noether currents follow.

\begin{proposition}[Noether and stress currents]
\label{prop:ms-Noether-currents}
Let $\mathcal L\mathcal F=0$ on $D$.

First, for any fixed $h$ satisfying $\mathcal Lh=0$, the vertical
transformation $\delta_h\mathcal F=h$ is a divergence variational symmetry.
Its Noether current is
\begin{equation}
\boldsymbol J_h
:=\frac{h\mathcal F_r-\mathcal Fh_r}{r}\,dz
-\frac{h\mathcal F_z-\mathcal Fh_z}{r}\,dr,
\qquad
d\boldsymbol J_h=0.
\label{eq:ms-Green-Noether-current}
\end{equation}

Second, translation in $z$ gives the divergence-free vector current
\begin{equation}
J_z^r=\frac{\mathcal F_r\mathcal F_z}{r},
\qquad
J_z^z=\frac{\mathcal F_z^2-\mathcal F_r^2}{2r},
\qquad
\partial_rJ_z^r+\partial_zJ_z^z=0.
\label{eq:ms-z-Noether-current}
\end{equation}

Third, define the weighted stress tensor
\begin{equation}
T_{ij}
:=\frac{\mathcal F_i\mathcal F_j}{r}
-\delta_{ij}\frac{|\nabla\mathcal F|^2}{2r}.
\label{eq:ms-stress-tensor}
\end{equation}
Then
\begin{equation}
\partial_iT_{iz}=0,
\qquad
\partial_iT_{ir}=\frac{|\nabla\mathcal F|^2}{2r^2}.
\label{eq:ms-stress-divergence}
\end{equation}
The bundle dilation $(r,z,\mathcal F)\mapsto
(e^s r,e^s z,e^{s/2}\mathcal F)$ gives, up to an overall sign, the improved
current
\begin{equation}
D_i:=x^jT_{ij}-\frac{\mathcal F\mathcal F_i}{2r},
\qquad
\partial_iD_i=0,
\qquad
(x^r,x^z)=(r,z).
\label{eq:ms-dilation-current}
\end{equation}
In particular, the unimproved stress identity gives the exact Pohozaev
formula
\begin{equation}
\frac12\int_D\frac{|\nabla\mathcal F|^2}{r}\,dr\,dz
=\int_{\partial D}x^jT_{ij}n_i\,ds.
\label{eq:ms-Pohozaev-identity}
\end{equation}
\end{proposition}

\begin{proof}
For the first statement,
\begin{equation}
\left.\frac{d}{d\epsilon}L_{\mathcal F+\epsilon h}
\right|_{\epsilon=0}
=\partial_i\!\left(\frac{\mathcal Fh_i}{r}\right)
\end{equation}
because $\mathcal Lh=0$. Equation~\eqref{eq:ms-Green-Noether-current} is
the resulting Noether current. The remaining identities follow by direct
differentiation and Eq.~\eqref{eq:ms-weighted-operator}. Integrating
$\partial_i(x^jT_{ij})=|\nabla\mathcal F|^2/(2r)$ proves
Eq.~\eqref{eq:ms-Pohozaev-identity}.
\end{proof}

Pure amplitude scaling $\mathcal F\mapsto a\mathcal F$ preserves the linear
field equation but multiplies the action by $a^2$; it is therefore a solution
symmetry, not a Noether symmetry. Translation in $r$ is broken by the
explicit weight $r^{-1}$. None of the currents above is a sign-definite
physical mass current.

\subsection{A formal Hamiltonian slice and its instability}

The covariant theory permits a canonical split, but no meridional coordinate
is physical time. To make this distinction explicit, let
$I=(r_-,r_+)\Subset(0,\infty)$ and use $z$ as a formal evolution variable.
Set
\begin{equation}
\pi(r):=\frac{\mathcal F_z(r,z)}{r},
\qquad
\Omega_I:=\int_I\delta\mathcal F\wedge\delta\pi\,dr,
\label{eq:ms-slice-symplectic-form}
\end{equation}
and
\begin{equation}
H_z[\mathcal F,\pi]
:=\int_I\left(\frac r2\pi^2-\frac{\mathcal F_r^2}{2r}\right)dr.
\label{eq:ms-z-Hamiltonian}
\end{equation}
For variations obeying boundary conditions that remove the endpoint term,
the Hamilton equations are
\begin{equation}
\mathcal F_z=\frac{\delta H_z}{\delta\pi}=r\pi,
\qquad
\pi_z=-\frac{\delta H_z}{\delta\mathcal F}
=-\partial_r\!\left(\frac{\mathcal F_r}{r}\right),
\label{eq:ms-z-Hamilton-equations}
\end{equation}
and hence reproduce $\mathcal L\mathcal F=0$. The corresponding formal
Poisson bracket is
\begin{equation}
\{P,Q\}_I
:=\int_I\left(
\frac{\delta P}{\delta\mathcal F}\frac{\delta Q}{\delta\pi}
-\frac{\delta P}{\delta\pi}\frac{\delta Q}{\delta\mathcal F}
\right)dr.
\label{eq:ms-slice-Poisson-bracket}
\end{equation}

For two solution translations $f$ and $g$, the bilinear form
\begin{equation}
c_I(f,g)
:=\int_I\frac{f g_z-g f_z}{r}\,dr
\label{eq:ms-solution-cocycle}
\end{equation}
is independent of $z$ when the radial endpoint flux vanishes. It is an
antisymmetric two-cocycle on the abelian solution-translation algebra. It
may be degenerate for a chosen boundary condition and is not a quantum
central extension.

The Hamiltonian in Eq.~\eqref{eq:ms-z-Hamiltonian} is indefinite. The exact
high-frequency Bessel solutions used in the Hadamard-instability theorem
therefore have a complementary interpretation: the formal $z$-Hamiltonian
flow amplifies arbitrarily small smooth Cauchy data exponentially. This is
an elliptic Cauchy evolution, not a unitary or causal dynamics. In
particular, no causal Peierls bracket follows from
Eq.~\eqref{eq:ms-slice-Poisson-bracket}.

\subsection{Boundary Lagrangians, gluing, and the physical interface}

For a fixed smooth boundary component $\Sigma$, define the conormal momentum
\begin{equation}
p_n:=\frac{\partial_n\mathcal F}{r}.
\label{eq:ms-conormal-momentum}
\end{equation}
The trace space
\begin{equation}
\mathcal H_\Sigma
:=H^{1/2}(\Sigma)\oplus H^{-1/2}(\Sigma)
\label{eq:ms-boundary-space}
\end{equation}
has the canonical symplectic form
\begin{equation}
\omega_\Sigma\bigl((f,p),(g,q)\bigr)
:=\int_\Sigma(fq-gp)\,ds.
\label{eq:ms-boundary-symplectic-form}
\end{equation}
Integration of Eq.~\eqref{eq:ms-form-formula} proves that the Cauchy data of
$\mathcal L$-harmonic fields are isotropic. Under the standard elliptic trace
and unique-continuation hypotheses, the Calder\'on projector identifies
their completed Cauchy-data space as a Lagrangian subspace of
$\mathcal H_{\partial D}$
\cite{BoossBavnbekLeschZhu2009JGP,CoxEtAl2016TAMS}. This functional-analytic
conclusion is stronger than the formal Green identity and uses the stated
trace and closed-range hypotheses.

On a fixed boundary and an invertible sheet,
\begin{equation}
\delta\mathcal F=A\,\delta\eta,
\qquad
\pi_n:=A p_n,
\qquad
\delta\mathcal F\wedge\delta p_n
=\delta\eta\wedge\delta\pi_n.
\label{eq:ms-boundary-canonical-map}
\end{equation}
Thus the auxiliary trace transformation is symplectic for $A\neq0$ and
collapses at $A=0$.

Suppose $D=D_-\cup_\Sigma D_+$ and orient both copies of $\Sigma$ by their
outward normals. The two variational boundary terms cancel precisely when
$\mathcal F$ and its conormal momentum match. Hence the solution spaces of
the two pieces define Lagrangian correspondences, and source-free gluing is
their composition along the diagonal. If $\mathcal F$ is continuous but
$p_n$ jumps, the uncancelled term is exactly the auxiliary delta-layer
source. This statement concerns the linear auxiliary equation: its jump is
not an Israel surface stress.

\begin{proposition}[Exact distributional interface dictionary]
\label{prop:ms-distributional-interface}
Let $\Sigma\Subset\{r>0\}$ be a smooth oriented curve with unit normal $n$
pointing from $D_-$ to $D_+$, and let $\mathcal F$ be piecewise $C^2$ with
one-sided traces. In the sense of Schwartz distributions,
\begin{align}
\mathcal L\mathcal F
={}&(\mathcal L\mathcal F)_{\rm reg}
+[\partial_n\mathcal F]\delta_\Sigma
+\operatorname{div}\!\left([\mathcal F]n\delta_\Sigma\right)
-\frac{n_r}{r}[\mathcal F]\delta_\Sigma.
\label{eq:ms-discontinuous-interface}
\end{align}
If $[\mathcal F]=0$, only the single-layer term remains. Within one
constant-$\ellzero$ chart, continuity of $\eta$ and $A\neq0$ give
\begin{equation}
[\partial_n\mathcal F]=A[\partial_n\eta].
\label{eq:ms-continuous-interface-eta}
\end{equation}
The exact reflected mode
\begin{equation}
\mathcal F(r,z)=rJ_1(kr)e^{-k|z|}
\label{eq:ms-reflected-layer-mode}
\end{equation}
satisfies
\begin{equation}
\mathcal L\mathcal F=-2krJ_1(kr)\delta(z).
\label{eq:ms-reflected-layer-source}
\end{equation}
It is a razor-thin auxiliary layer, not a smooth volume-dust solution.
\end{proposition}

\begin{proof}
Write the piecewise field with a Heaviside function of signed distance to
$\Sigma$ and differentiate distributionally. The regular gradients produce
$[\partial_n\mathcal F]\delta_\Sigma$, the jump of the field produces the
divergence term, and $-r^{-1}\partial_r$ produces the last term in
Eq.~\eqref{eq:ms-discontinuous-interface}. If the field is continuous, the
last two terms vanish. On one inverse chart,
$\partial_n\mathcal F=B n_r+A\partial_n\eta$; continuity of $r$, $\eta$,
$A$, and $B$ proves Eq.~\eqref{eq:ms-continuous-interface-eta}. Finally,
$\partial_z^2e^{-k|z|}=k^2e^{-k|z|}-2k\delta(z)$, while the Bessel equation
gives
$$
\left(\partial_r^2-r^{-1}\partial_r+k^2\right)
\bigl[rJ_1(kr)\bigr]=0,
$$
which proves Eq.~\eqref{eq:ms-reflected-layer-source}.
\end{proof}

The wavefront of $\delta_\Sigma$ is the nonzero conormal bundle of
$\Sigma$. Ellipticity therefore localizes the singular support of the
auxiliary response to that of its source, rather than producing propagation
along real characteristics \cite{Hormander2003ALPDOI}. If $\eta$ or the
metric itself jumps, expressions such as $|\nabla\eta|^2$, $\log\eta$,
$g^{-1}$, and the nonlinear Einstein tensor contain products that ordinary
distribution theory need not define. They must not be formed without a
separately specified generalized-function framework or regularization.
Moreover, the coefficient in Eq.~\eqref{eq:ms-discontinuous-interface} is
not an Israel stress tensor: that tensor is determined only after the full
metric has been reconstructed and its extrinsic-curvature jump evaluated.
The Geroch--Traschen regular-metric class accommodates hypersurface shells
but not generic codimension-two strings \cite{GerochTraschen1987PRD}.

For a fixed piecewise smooth boundary, the first-order scalar action creates
no additional codimension-two charge. Shape variations, moving interfaces,
or independent boundary actions enlarge the phase space and can generate
corner terms. A Maslov index or spectral-flow statement requires a specified
self-adjoint Fredholm family; it is not a consequence of the word
``Lagrangian'' alone.

The vacuum exterior has a different variational system. Where a global twist
potential $\chi$ exists, the gauge-fixed Ernst subsystem may be written, up
to an overall conventional constant, as
\begin{equation}
\mathcal A_{\mathrm E}[f,\chi]
=\frac12\int_D\frac r{f^2}
\left(|\nabla f|^2+|\nabla\chi|^2\right)dr\,dz,
\qquad f>0.
\label{eq:ms-Ernst-action}
\end{equation}
It has polymomenta
\begin{equation}
p_f^i=\frac r{f^2}\partial_i f,
\qquad
p_\chi^i=\frac r{f^2}\partial_i\chi,
\label{eq:ms-Ernst-polymomenta}
\end{equation}
and an analogous conserved multisymplectic current for pairs of Jacobi
fields. This structure belongs to the nonlinear Ernst harmonic map, not to
$\mathcal F$. On a multiply connected exterior, Eq.~\eqref{eq:ms-Ernst-action}
is global only after the twist periods vanish; otherwise one must use local
potentials or the universal cover together with the cohomology sector.

Finally, the physical variational problem starts from the
Einstein--Hilbert action with its gravitational boundary terms and a dust
action capable of vorticity. For example, a material-coordinate
representation uses
\begin{equation}
\mathcal A_{\mathrm{dust}}
=-\frac12\int_{\mathcal M}\sqrt{-g}\,M
\bigl(g^{ab}U_aU_b+1\bigr)d^4x,
\qquad
U_a=-\partial_aT+W_A\partial_aZ^A.
\label{eq:ms-vortical-dust-action}
\end{equation}
Such matter actions and the required gravitational boundary term are
reviewed in Refs.~\cite{Brown1993CQG,GibbonsHawking1977PRD}.
The full covariant presymplectic form has diffeomorphism degeneracies and
must be reduced only after the constraints, asymptotic conditions, and
boundary gauge have been fixed \cite{LeeWald1990JMP,IyerWald1994PRD}.
Variation at an internal timelike boundary
relates continuity of the induced metric and gravitational canonical
momentum to the Darmois conditions; the matter variables additionally
require a no-flux or prescribed-flux condition. No equality between that
physical form and $\Omega_{\mathcal F}$ has been established here.

This separation prevents three incorrect conclusions. The convex auxiliary
energy is not the ADM mass, the formal $z$-Hamiltonian is not physical time
evolution, and the auxiliary Lagrangian matching relation is only one
compatibility layer within the complete dust--vacuum boundary problem.
